% \fancyhf{} 
% \fancyfoot[C]{\thepage}


\chapter{PENDAHULUAN}
% \thispagestyle{plain} % Halaman pertama bab menggunakan gaya plain


\section{Latar Belakang}

\par Dalam era digital saat ini, banyak institusi pendidikan telah beralih ke platform pembelajaran berbasis web untuk menyediakan layanan pendidikan jarak jauh. Namun, tantangan utama yang sering dihadapi adalah kurangnya kemampuan sistem untuk memahami pola perilaku pengguna dan memberikan rekomendasi konten yang tepat sasaran \citep{rabiah2019manajemen}.

\par Algoritma K-Nearest Neighbors (KNN) merupakan salah satu metode \textit{supervised learning} yang cukup sederhana namun efektif dalam masalah klasifikasi dan regresi \citep{enwiki:1285981318}. Penelitian ini akan mengimplementasikan KNN untuk memprediksi kualitas pengguna berdasarkan data aktivitas pengguna pada aplikasi e-learning \citep{guo2003knn}.

% jika ingin mengutip langsung bisa pakai \cite
\par Menurut \cite{nichols2003theory}, Metode e-learning telah merevolusi cara individu mengakses pendidikan dan pelatihan, menawarkan fleksibilitas dan aksesibilitas yang tak tertandingi. Seiring dengan peningkatan adopsi platform e-learning berbasis web, memahami dan memprediksi kualitas pengalaman pengguna menjadi krusial untuk meningkatkan efektivitas pembelajaran. Kualitas pengguna yang tinggi tidak hanya berkorelasi dengan retensi dan keterlibatan, tetapi juga secara langsung memengaruhi keberhasilan tujuan pendidikan. Penelitian ini akan berfokus pada penerapan algoritma \textit{K-Nearest Neighbors} (KNN) untuk memprediksi kualitas pengguna pada aplikasi e-learning berbasis web, dengan mempertimbangkan berbagai faktor interaksi dan performa yang dapat memengaruhi pengalaman belajar daring.


\section{Rumusan Masalah}
Berdasarkan latar belakang yang telah diuraikan, permasalahan dalam penelitian ini dapat dirumuskan sebagai berikut:
\begin{enumerate}
    \item   Diperlukan implementasi algoritma KNN dalam memprediksi kualitas pengguna.
    \item Diperlukan analisa tingkat akurasi model dalam mengklasifikasikan kualitas pengguna
    \item Diperlukan umpan balik dari pengembang dan pengguna guna meningkatkan pengalaman belajar pada aplikasi e-learning berbasis web?
\end{enumerate}
\section{Tujuan Penelitian}  
Adapun tujuan dari penelitian ini adalah sebagai berikut: 
\begin{enumerate}
    \item MengimplementUntuk mengimplementasikan algoritma KNN dalam prediksi kualitas pengguna.
    \item  Untuk mengevaluasi tingkat akurasi model dalam memprediksi kualitas pengguna\textit{framework} MediaPipe untuk mengenali abjad BISINDO.
    \item Perbaikan pada sistem dan pelayanan berdasarkan umpan balik.
\end{enumerate}

%% subbagian tambahan optional
\section{Batasan permasalahan}

\par Penelitian ini dibatasi pada:
\begin{enumerate}
    \item Dataset dummy dengan jumlah 100 sample data.
    \item Fitur yang digunakan hanya waktu akses rata-rata, jumlah login, durasi video, dan rating pengguna.
    \item Hanya menggunakan algoritma KNN tanpa komparasi dengan algoritma lain.
\end{enumerate}

\section{Manfaat Penelitian}  
\par Adapun manfaat dari penelitian ini adalah sebagai berikut:  
\begin{enumerate}
    \item Meningkatkan kualitas layanan e-learning melalui personalisasi berbasis prediksi kualitas pengguna.
    \item Menjadi referensi bagi pengembangan sistem rekomendasi berbasis \textit{machine learning}.
    \item Perbaikan sistem berdasarkan umpan balik pada website
\end{enumerate}

% Baris ini digunakan untuk membantu dalam melakukan sitasi
% Karena diapit dengan comment, maka baris ini akan diabaikan
% oleh compiler LaTeX.
\begin{comment}
\bibliography{daftar-pustaka}
\end{comment}